\chapter{Subdural haematoma}
\index{Subdural haematoma}

\medskip
\noindent\textit{Educational reference; always seek professional advice for individual care.}

\section*{Definition and Overview}
A subdural haematoma is a collection of blood that accumulates between the dura mater, the tough outer membrane covering the brain, and the arachnoid mater, the delicate membrane beneath it. It is a form of intracranial haemorrhage caused by the tearing of bridging veins that cross this space. Because the skull is a rigid container, any expanding collection of blood compresses the brain, and the resulting pressure can cause neurological damage. The severity ranges widely: some subdural haematomas are small, produce few symptoms, and resolve on their own, while others expand rapidly and are life-threatening emergencies requiring urgent surgery.

Subdural haematomas are classically divided by timing into acute, subacute, and chronic forms, based on how quickly symptoms develop after the injury. Acute subdural haematomas become symptomatic within hours to days of a head injury and are often severe. Chronic subdural haematomas develop weeks to months after an injury, which may have been trivial and is sometimes not even remembered; they are seen mainly in older adults and in people taking blood-thinning medicines. The distinction matters because the mechanism, the population affected, the clinical picture, and the treatment differ.

This chapter explains why subdural haematomas occur, who is most at risk, how they present, and how they are diagnosed and treated. It also covers the complications to watch for, the outlook after treatment, and the practical steps that reduce the risk of developing one. The focus throughout is on recognising the condition early, because prompt assessment and treatment are the strongest predictors of a good outcome. Subdural haematoma is a condition where awareness genuinely saves lives, particularly among older people, who carry a disproportionately high share of the risk.

\section*{Epidemiology}
Subdural haematoma is one of the commoner serious consequences of head injury. Acute subdural haematoma complicates a small but important proportion of significant head injuries, and it is the most dangerous of the surgically treatable intracranial haemorrhages. Severe acute cases are most often seen in younger adults, particularly after motor vehicle crashes, falls from a height, and assaults, and in these settings the injury frequently involves more than just the bleeding. A large proportion of these patients have an associated injury to the brain tissue itself, which worsens the outlook.

Chronic subdural haematoma has a quite different pattern. It is predominantly a condition of older age, becoming increasingly common after the sixth decade of life, and its incidence has been rising in many countries and other high-income countries. This rise is linked to an ageing population and to the widespread use of anticoagulant medicines, which make bleeding more likely after even minor head trauma. Men are affected about two to three times as often as women. Additional risk factors include alcohol misuse, epilepsy, and a history of falls. Chronic subdural haematoma is now common enough that it should be considered in any older person who develops confusion, drowsiness, or a change in function over days to weeks, especially with any recent history of a fall, however minor.

\section*{Aetiology and Pathophysiology}
The bridging veins that drain blood from the surface of the brain into the large venous sinuses cross the subdural space and are stretched across it. When the head is subjected to acceleration and deceleration, the brain moves within the skull, and these veins can tear at the point where they attach to the sinus. The result is bleeding into the subdural space.

In acute subdural haematoma, the bleeding is brisk and the collection forms quickly. As the blood accumulates, it raises pressure inside the skull, which impairs blood flow to the brain and can shift the brain structures sideways or downward, compressing vital regions such as the brainstem. In chronic subdural haematoma, the initial bleed is often small and stops, but the clot then breaks down, attracting an inflammatory reaction and the growth of fragile new blood vessels that ooze repeatedly, gradually enlarging the collection over weeks.

\begin{itemize}
    \item \textbf{Head injury:} the commonest cause; the injury may be severe, as in a high-speed crash, or minor, as in a stumble against furniture
    \item \textbf{Anticoagulant and antiplatelet medicines:} warfarin, other blood thinners, and aspirin increase the risk of bleeding and make even trivial trauma dangerous
    \item \textbf{Older age:} the brain shrinks with age, which stretches the bridging veins, and the supporting tissue becomes more fragile
    \item \textbf{Alcohol misuse:} contributes through falls, impaired clotting, and brain shrinkage
    \item \textbf{Cerebrospinal fluid drainage:} in people with shunts or after brain surgery, the draining reduces brain volume and stretches the bridging veins
    \item \textbf{Blood disorders:} conditions that impair clotting, such as haemophilia or low platelet counts, predispose to bleeding
\end{itemize}

\section*{Clinical Features}
The presentation depends on how quickly the haematoma enlarges and how much the brain is compressed.

\begin{itemize}
    \item \textbf{Headache:} often persistent and worsening, and a prominent feature especially in chronic cases
    \item \textbf{Confusion and drowsiness:} a fluctuating level of alertness, increasingly recognised as a common presentation in older people
    \item \textbf{Weakness of one side:} arm or leg weakness, typically on one side of the body, reflecting pressure on the motor pathways
    \item \textbf{Personality and memory changes:} irritability, apathy, or impaired memory that may be mistaken for dementia in older patients
    \item \textbf{Slurred speech:} difficulty speaking clearly, sometimes with difficulty understanding others
    \item \textbf{Seizures:} a convulsion may be the first obvious event, especially in acute cases
    \item \textbf{Nausea and vomiting:} common in acute bleeds because of raised pressure inside the skull
    \item \textbf{Unequal pupils:} one pupil larger than the other, a late and serious sign indicating compression of the eye-movement nerve
    \item \textbf{Loss of consciousness:} in severe acute cases the person may never regain consciousness after the injury
\end{itemize}

Any combination of these features following a head injury, or developing gradually in an older person, warrants urgent medical assessment.

\section*{Differential Diagnosis}
Several other conditions can look like a subdural haematoma, and imaging is usually needed to tell them apart.

\begin{itemize}
    \item \textbf{Epidural haematoma:} bleeding between the skull and the dura, usually arterial and faster, with a characteristic lucid interval before rapid deterioration; requires urgent surgery
    \item \textbf{Subarachnoid haemorrhage:} bleeding into the fluid-filled space around the brain, typically a sudden severe "thunderclap" headache, often from a ruptured aneurysm
    \item \textbf{Intracerebral haemorrhage:} bleeding within the brain tissue itself, producing focal deficits such as weakness or speech loss
    \item \textbf{Ischaemic stroke:} a blocked blood vessel rather than bleeding; may be indistinguishable without a scan and is treated differently
    \item \textbf{Dementia:} gradual memory and behaviour change in older people can mimic the subacute cognitive decline of a chronic subdural haematoma
    \item \textbf{Migraine:} recurrent headache with nausea, but without the focal neurological signs, drowsiness, or recent trauma
    \item \textbf{Brain tumour:} can cause progressive headache, seizures, and focal signs over weeks
\end{itemize}

\section*{When to Seek Medical Care}
Anyone with a head injury, especially if they take blood-thinning medicines, have risk factors, or are older, should be assessed, and a subdural haematoma should be suspected if there is any persistent headache, confusion, or drowsiness. A person who develops personality change, poor balance, or difficulty with memory over weeks should also be reviewed, even without a remembered injury.

\textbf{Contact your local emergency services for:}
\begin{itemize}
    \item Loss of consciousness or a reduced level of alertness at any time after a head injury
    \item Worsening headache, repeated vomiting, or a seizure after a head injury
    \item New weakness of an arm or leg, difficulty speaking, or confusion
    \item Unequal pupils, a severe fall in consciousness, or inability to wake the person
\end{itemize}

\section*{Diagnosis and Investigations}
\begin{itemize}
    \item \textbf{Computed tomography (CT) scan:} the first-line test; it shows blood clearly and quickly and is available in all local emergency departments
    \item \textbf{Magnetic resonance imaging (MRI):} more sensitive for smaller or older collections and for damage to the brain tissue itself; used when CT is inconclusive or as follow-up
    \item \textbf{History and examination:} the circumstances of any injury, medicines taken, and a neurological examination guide the urgency and the interpretation of scans
    \item \textbf{Blood tests:} clotting studies and a full blood count check for bleeding disorders and are essential before any surgery
    \item \textbf{Repeated scans:} chronic subdural haematomas can enlarge, so follow-up imaging may be arranged to monitor progress
    \item \textbf{Specialist review:} a neurosurgeon assesses whether surgery is needed and when
\end{itemize}

\section*{Management}
\begin{itemize}
    \item \textbf{Observation and monitoring:} small haematomas in stable patients may be managed with close neurological observation in hospital and repeat imaging
    \item \textbf{Reversing anticoagulation:} urgent reversal of warfarin or other blood thinners is often the first step, using specific antidotes or clotting factors
    \item \textbf{Burr hole drainage:} the standard surgical treatment for chronic subdural haematomas, in which a small hole is drilled in the skull and the blood is drained; it is quick, effective, and well tolerated even in frail patients
    \item \textbf{Craniotomy:} for acute haematomas and clots that are solid, a larger opening in the skull allows the clot to be removed and bleeding points to be controlled
    \item \textbf{Managing brain pressure:} medicines and positioning that reduce swelling and pressure inside the skull are used while recovery proceeds
    \item \textbf{Seizure treatment:} antiepileptic medicines are given if seizures occur
    \item \textbf{Rehabilitation:} physiotherapy, occupational therapy, and speech therapy help restore strength, balance, and function after discharge
    \item \textbf{Treatment review:} decisions about restarting blood thinners are made carefully with the treating team, balancing the risk of another bleed against the reason the medicine was prescribed
\end{itemize}

\section*{Complications}
\begin{itemize}
    \item \textbf{Brain compression:} a rapidly expanding haematoma can push the brain against rigid structures, leading to brainstem injury and death if untreated
    \item \textbf{Recurrence:} chronic subdural haematomas recur after drainage in a significant minority of cases, sometimes requiring repeat procedures
    \item \textbf{Postoperative infection:} bleeding into a drained cavity, infection, or new bleeding can complicate surgery
    \item \textbf{Seizures:} epilepsy may develop after the injury or the surgery
    \item \textbf{Neurological disability:} weakness, memory problems, or personality change can persist, especially after severe acute bleeds with associated brain injury
    \item \textbf{Secondary brain damage:} reduced blood flow and oxygen to compressed brain tissue adds to the injury from the original trauma
\end{itemize}

\section*{Prognosis and Outcomes}
The outcome of a subdural haematoma depends strongly on its type and on the speed of treatment. Chronic subdural haematomas, the form seen mainly in older adults, respond well to drainage: most people improve substantially, often within days, and the majority regain their previous level of function. Recurrence is common enough that follow-up imaging and a period of activity restriction are routine, but repeat drainage usually settles the problem. The outlook for acute subdural haematoma is more guarded and depends chiefly on the level of consciousness at presentation and on whether the brain itself was injured. Patients who are alert and are treated promptly generally do well, while those who present deeply unconscious have a much poorer prognosis despite surgery. With modern neurocritical care, rehabilitation, and coordinated discharge planning, many people make meaningful recoveries, and every month of improvement is valuable.

\section*{Prevention and Self-Care}
\begin{itemize}
    \item Prevent falls by keeping walkways clear, using handrails, ensuring good lighting, and reviewing medicines that cause dizziness with the doctor
    \item Use seatbelts and appropriate child restraints in motor vehicles, and wear a helmet when cycling or riding a motorcycle
    \item If you take blood-thinning medicines, keep all appointments for monitoring, report any head injury promptly, and carry a card listing your medicines
    \item Limit alcohol consumption, which raises the risk of falls and of bleeding
    \item Wear a helmet for activities with a high fall risk, such as horse riding, skiing, and climbing ladders
    \item After any head injury, even a minor one, watch for headache, drowsiness, confusion, vomiting, or weakness, and seek medical advice if any develop
    \item Follow up after discharge from hospital, including attending repeat scans and rehabilitation sessions as arranged
\end{itemize}

\section*{Sources and Further Reading}
The following organisations provide reliable information about head injuries and intracranial bleeding for patients and families.

\begin{itemize}
    \item NHS. \textit{Head injury and subdural haematoma information.} https://www.nhs.uk/
    \item Mayo Clinic. \textit{Subdural haematoma: symptoms, causes, and treatment.} https://www.mayoclinic.org/
    \item healthdirect Australia. \textit{Head injury and intracranial bleeding.} https://www.healthdirect.gov.au/
    \item World Health Organization (WHO). \textit{Injuries and violence prevention.} https://www.who.int/
    \item Australian Department of Health and Aged Care. \textit{Injury prevention and management resources.} https://www.health.gov.au/
    \item Headway Australia. \textit{Support and information after acquired brain injury.} https://headwayaustralia.org.au/
\end{itemize}

\clearpage
